% GPSS 1980, Lemma 1 and positive decomposition preparation, pp. 169--171.

\subsection{The GPSS special normal form and positive decomposition construction}
\label{sec:p09-gpss-special-normal-form}

This presentation supplies the first-order normal form and the positive
construction steps in the GPSS expressive-completeness outline. The
normal form works on every nonempty strict linear order. The subsequent
negation proposition is a separate obligation in the source's proof for
discrete complete orders
\cite[pp.~169--171]{GabbayPnueliShelahStavi1980}.

\begin{hbdefinition}{def:p09-gpss-regional-stages}{The literal GPSS atoms and signed regional stages}
The first-order language \(L_1\) has equality, strict order \(<\),
unary predicates \(q_i\), and the classical connectives and quantifiers.
Its models have nonempty domain \(T\), with arbitrary unary
interpretations. The subscript \(G\) below distinguishes the paper's
classes from Kamp's different classes. Put
\[
              \operatorname{At}_G(z)=\{q_i(z),\neg q_i(z):i\geq0\}.
\]
These are individual signed literals, not arbitrary Boolean formulas
or arbitrary one-variable first-order conditions.

An interval endpoint may be a variable, the formal lower endpoint
\(-\infty\), or the formal upper endpoint \(+\infty\), in the
appropriate position. The formal endpoints are never quantified or
interpreted as elements of \(T\). Let \(g_{ab}(z)\) be
\(a<z\land z<b\), omitting a comparison when its endpoint is
formal. Thus \(g_{-\infty,+\infty}(z)\) is true. Define
\[
 \operatorname{Bet}_{G,0}(a,b)=\varnothing,
\]
and let \(\operatorname{Bet}_{G,n+1}(a,b)\) consist of
\[
 e=\exists z\left(g_{ab}(z)\land\bigwedge_{j=1}^{m}\gamma_j\right)
                  \quad\text{and}\quad\neg e,
\]
where \(m\geq0\), the bound variable is fresh, and every
\(\gamma_j\) belongs to
\[
 \operatorname{Bet}_{G,n}(a,z)\cup\operatorname{At}_G(z)
                                  \cup\operatorname{Bet}_{G,n}(z,b).
\]
An empty conjunction is true. Each index is a nonnegative integer.
For actual endpoint variables use distinct names; their assigned values
need not be distinct unless an order prefix requires it. Write
\[
\begin{aligned}
 \operatorname{Bet}_G(a,b)&=\bigcup_{n\geq0}\operatorname{Bet}_{G,n}(a,b),\\
 \operatorname{Bef}_{G,n}(x)&=\operatorname{Bet}_{G,n}(-\infty,x),\qquad
 \operatorname{Aft}_{G,n}(x)=\operatorname{Bet}_{G,n}(x,+\infty).
\end{aligned}
\]
Unindexed ray classes are their stage unions. Using both formal endpoints
also gives whole-order components, needed below for sentences. This makes
the zero-free-variable extension of the source convention explicit
\cite[p.~170]{GabbayPnueliShelahStavi1980}.
\end{hbdefinition}

\begin{hbdefinition}{def:p09-gpss-special-form}{Special clauses with complete weak order types}
For a finite nonempty list \(X\) of distinct parameter variables,
a complete weak order type is an ordered partition
\((B_1,\ldots,B_s)\) of \(X\). Choose a representative
\(y_j\in B_j\). Its prefix \(\pi\) conjoins
\(y_1<\cdots<y_s\) and the equalities of every variable
in \(B_j\) with \(y_j\). For one representative use its
self-equality when the prefix otherwise has no conjunct.

A GPSS special clause on this type is
\(\pi\land\bigwedge_j\gamma_j\), where its local conjuncts
come from
\[
 \operatorname{Bef}_G(y_1)\cup
 \bigcup_{j=1}^{s}\operatorname{At}_G(y_j)\cup
 \bigcup_{j=1}^{s-1}\operatorname{Bet}_G(y_j,y_{j+1})\cup
 \operatorname{Aft}_G(y_s).
\]
The list of local conjuncts may be empty. A special normal form is a
finite disjunction of such clauses; the empty disjunction is false.
Coincident parameters form one block and do not create a spurious
strict interval. For an empty parameter list, a special sentence
clause is a finite conjunction of whole-order components, with no
prefix; finite disjunctions and their empty cases have the same meanings.
All substitutions concern free occurrences and avoid variable capture
\cite[Lemma~1, p.~170]{GabbayPnueliShelahStavi1980}.
\end{hbdefinition}

\begin{hblemma}{thm:p09-gpss-kamp-normal-form-bridge}{proved=annex:p9.gpss.kamp.normal.form.bridge}{The exact bridge to Kamp's earlier regional grammar}
The GPSS regional stages are increasing, and signed atoms and regional
components have a complement operation given by reversing their sign.
Finite Boolean combinations of them can therefore be put effectively
into finite disjunctive normal form using the same permitted components.

After renaming \(q_i\) as \(P_i\), replacing equality by
\(\neg(x<y)\land\neg(y<x)\), and using the starred
Kamp classes of \cref{def:p09-kamp-iii1-components}, each GPSS
before, after or between component with at least one actual endpoint
variable translates to a corresponding
Kamp component. Conversely each starred Kamp component translates
to a finite disjunction of conjunctions of GPSS components in the
same region. A Kamp pointwise literal conjunction becomes a list
of GPSS point literals; its distinguished true and false slots
become respectively an empty conjunction and empty disjunction.

Both translations preserve meaning on every nonempty strict linear
order and introduce no new unary predicate. They provide an explicit
bridge to \cref{thm:p09-kamp-iii1}; they do not identify the two
syntactic classes. In particular Kamp's later arbitrary elementary
one-point formulas are not GPSS \(\operatorname{At}_G\) literals.
Whole-order components use the separate sentence convention.
\end{hblemma}

\begin{hbtheorem}{thm:p09-gpss-special-normal-form}{proved=annex:p9.gpss.special.normal.form}{The complete GPSS special-form normal-form lemma}
Every \(L_1\)-formula is effectively equivalent, on every
nonempty strict linear order and every unary interpretation, to a
finite disjunction of GPSS special clauses on its declared parameter
list. The construction retains all weak order types, including
equal-point blocks, and handles every existential and universal
quantifier. Sentences have the explicit whole-order version above.
No unary predicate absent from the input is needed.

The normal form still contains nested first-order regional
quantification. This result is GPSS's Lemma~1, with its induction,
equality cases, sentence convention and empty-list cases supplied;
it is not yet the later conversion to future temporal syntax
\cite[Lemma~1, p.~170]{GabbayPnueliShelahStavi1980}.
\end{hbtheorem}

\begin{hbcorollary}{thm:p09-gpss-future-special-form}{proved=annex:p9.gpss.future.special.form}{The exact present-and-future specialization}
Suppose \(\varphi(x)\) is a source future formula: every
quantifier is restricted to \(y\geq x\), relative to the
same distinguished free variable. Then it is equivalent on every
linear-order model to a finite disjunction of conjunctions from
\(\operatorname{At}_G(x)\cup\operatorname{Aft}_G(x)\).
The inclusive lower bound retains the present atomic conditions.
No before component, past-dependent one-point label, or completeness
hypothesis is needed for this first-order specialization
\cite[pp.~169--170]{GabbayPnueliShelahStavi1980}.
\end{hbcorollary}

\subsubsection{The exact future-labelled decomposition class}

\begin{hbdefinition}{def:p09-gpss-future-labels}{Permitted labels and their strict future semantics}
A permitted label is the first-order translation \(B^*(u)\)
of a finite \(L(U)\)-formula \(B\), with no free individual
parameter except its current point \(u\). Its strict until clause is
\[
 (B\mathbin{\mathsf U_s}C)^*(u)=
 \exists v>u\left(C^*(v)\land
                  \forall w(u<w<v\to B^*(w))\right).
\]
Atomic and Boolean clauses are homomorphic. Put
\(FB=\top\mathbin{\mathsf U_s}B\) and \(GB=\neg F\neg B\).
True and false are ordinary Boolean abbreviations, for example
\(q_0\lor\neg q_0\) and \(q_0\land\neg q_0\), with
no semantic dependence on that inessential letter.
Before copying a label to a chain point, rename its bound variables
apart from every endpoint and chain variable, and replace only its
free current-point variable.
All labels are evaluated in the original whole model. A label at an
interior point may inspect points beyond the displayed interval's
right endpoint; it is not implicitly relativized to that interval
\cite[pp.~164, 169--170]{GabbayPnueliShelahStavi1980}.
\end{hbdefinition}

\begin{hblemma}{thm:p09-gpss-future-label-closure}{proved=annex:p9.gpss.future.label.closure}{Boolean closure and the exact future-only restriction}
Permitted labels are closed under all finite Boolean operations.
Their truth at a point depends only on the ordered labelled tail
at and after that point. Thus a general past-dependent formula
cannot be inserted as a decomposition label merely because it has
one free variable. In particular \(\exists v<u\,q_0(v)\)
has no equivalent future label over all linear-order models.
\end{hblemma}

\begin{hbdefinition}{def:p09-gpss-decomposition}{Decompositions with finite or formal endpoints}
For admissibly ordered endpoints \(a<b\), choose \(r\geq0\)
interior points \(z_1<\cdots<z_r\) of \(T\), all between the
endpoints. There are \(r\) point labels \(\theta_i\), and
\(r+1\) gap labels \(\phi_i\), all permitted future labels.
Put \(z_0=a,z_{r+1}=b\) only as endpoint abbreviations. The
decomposition formula is
\[
\begin{aligned}
 \delta(a,b)={}&a<b\ \land\
 \exists z_1\cdots z_r\bigl[
       a<z_1<\cdots<z_r<b\\
 &\qquad\land\bigwedge_{i=1}^{r}\theta_i(z_i)
 \land\bigwedge_{i=0}^{r}
       \forall u(z_i<u<z_{i+1}\to\phi_i(u))\bigr].
\end{aligned}
\]
Formal endpoint comparisons and quantification are omitted as in the
regional grammar. For \(r=0\), the internal chain and point-label
conjunction are empty and there is just the one gap universal.
In the paper's finite-endpoint notation this is \(n=r+1\),
with \(z_0=x,z_n=y\). Endpoints have no point labels. All
open gaps may be empty even when their endpoints are distinct.
The finite-endpoint spelling with \(n=0\) would require
\(x=y\) as well as \(x<y\), and contributes only false.

If the right endpoint is \(+\infty\), also permit a finite
list of global cofinal-occurrence conditions
\[
          \operatorname{Inf}(\psi)
             :=\forall u\in T\ \exists v>u\ \psi(v),
\]
where \(\psi\) is a permitted future label. The quantifier
over \(u\) is over the whole domain. Empty lists impose no
condition. Finite disjunctions of decompositions, including the
empty disjunction, will be used. Statements about complements on
a finite interval are relative to its guard \(a<b\)
\cite[pp.~170--171]{GabbayPnueliShelahStavi1980}.
\end{hbdefinition}

\begin{hbtheorem}{thm:p09-gpss-positive-decomposition}{proved=annex:p9.gpss.positive.decomposition}{Finite positive closure and existential assembly}
Finite disjunctions of the stated decompositions are closed under
finite conjunction on a fixed admissible interval. The conjunction
construction enumerates all order-preserving merges of the two
interior witness chains, including coincident points from different
chains. It conjoins point labels and transfers each covering
gap label onto newly inserted points and subintervals. Every resulting
label remains future-only. Cofinal-condition lists are combined by
union when the upper endpoint is formal.

If \(L(a,z)\) and \(R(z,b)\) have such finite disjunctions,
and \(\theta(z)\) is a permitted future label, then so does
\[
 \exists z\bigl(g_{ab}(z)\land L(a,z)\land\theta(z)\land R(z,b)\bigr).
\]
The construction concatenates the two witness chains and puts
\(\theta\) at their common interior point. It includes both rays
and the whole-order endpoint convention.

The empty conjunction is represented, on the admissible interval,
by the decomposition with no interior point and gap label true;
the empty disjunction is false. On admissible endpoints, every member of
\(\operatorname{Bet}_{G,1}\), of either sign, already has a
decomposition presentation. More generally the positive regional
constructor is handled whenever the presentations of all its
regional children are supplied. No arbitrary negation closure is
assumed in these claims
\cite[pp.~170--171]{GabbayPnueliShelahStavi1980}.
\end{hbtheorem}

\begin{hblemma}{thm:p09-gpss-suffix-translation}{proved=annex:p9.gpss.suffix.translation}{The complete translation of a future decomposition}
For a decomposition \(\delta(x,+\infty)\) with actual lower
endpoint \(x\), let \(\Theta_i,\Phi_i,\Psi_j\) be
the \(L(U)\)-formulas supplying its labels. Define
\[
 H_r=G\Phi_r,\qquad
 H_i=\Phi_i\mathbin{\mathsf U_s}(\Theta_{i+1}\land H_{i+1})
                    \quad(0\leq i<r).
\]
Then the exact future expression is
\[
                 H_0\land\bigwedge_j(F\Psi_j\land GF\Psi_j).
\]
It is equivalent to the decomposition at every point of every
nonempty linear-order model, including a last point. For \(r=0\)
the first factor is simply \(G\Phi_0\). This proves the
source's positive return from future decompositions to \(L(U)\)
\cite[p.~171]{GabbayPnueliShelahStavi1980}.
\end{hblemma}

\begin{hbexample}{def:p09-gpss-positive-merge-example}{Equal witnesses and the labels at inserted points}
Let \(p,q,r,s\) be proposition letters. A decomposition with
one interior point labelled \(p\) and both gaps labelled true
expresses \(D_p(x,y)=\exists z(x<z<y\land p(z))\) on
\(x<y\). The conjunction \(D_p\land D_q\) has three
decomposition cases: one point labelled \(p\land q\); two
points labelled \(p,q\) in that order; or two points labelled
\(q,p\). All their gap labels are true. On the three-point
order \(\{0,1,2\}\), with endpoints zero and two and both
letters true at one, only the coincident-witness case can hold.

If instead the first decomposition has point label \(p\)
and gap labels \(r,s\), and the other has no interior point
and gap label \(q\), their intersection has point label
\(p\land q\) and gap labels \(r\land q,s\land q\).
The \(q\)-condition at the merged point is necessary: in
the same three-point order the two smaller open gaps are empty,
but the original single gap still contains the middle point.
\end{hbexample}

\begin{proofboundary}{def:p09-gpss-negation-boundary}{The exact endpoint scope of complementation}
On discrete complete orders, an actual lower endpoint permits the
complete relative-complement constructions of
\cref{thm:p09-gpss-bounded-complement,thm:p09-gpss-suffix-complement}.
For finite endpoints the target is
\(x<y\land\neg\delta(x,y)\); for a future ray it is
\(\neg\delta(x,+\infty)\).
The next presentation supplies the event grammar, determining recursion
and complete translation of its defined, undefined and weak-order cases.
These results close the regional induction and yield
\cref{thm:p09-gpss-discrete-future-completeness}.

The unchanged future-label class is not closed under complementation
at a formal left endpoint. The prefix and whole-order counterexamples
are proved in \cref{thm:p09-gpss-prefix-complement-obstruction}.
Arbitrary complete orders require their separate nonattainment and
cut constructions
\cite[p.~171]{GabbayPnueliShelahStavi1980}.
\end{proofboundary}
